\subsection{Baselines}\label{apdx:baselines}

We compare the discrete- and continuous-time models against a set of baselines that isolate
the contribution of (i) class imbalance, (ii) temporal persistence, (iii) intrinsic agent
preferences, and (iv) social influence. Writing $d_p$ and $d_q$ for the persona- and
question-embedding dimensions, we report the parameter count of each so that the comparison in
Table~\ref{tab:prediction_models} can be read against model capacity. All baselines with
parameters are fit on the same transitions and with the same objective as the update rules
(Appendix~\ref{apdx:fitting}).

\textit{Majority Class.} Majority class is the simplest possible predictor. It ignores agent identity and interactions and always predicts the majority opinion observed in the training data,
\[
\hat{s}_i(t{+}1)=c,
\qquad
c=\operatorname{sign}\!\Big(\textstyle\sum_{\text{train}} s\Big).
\]
The model contains a single learned parameter, the constant label $c\in\{+1,-1\}$, and measures
the extent to which performance can be explained by class imbalance alone. We report this baseline only in tables that report accuracy (not balanced) because some models exhibit label bias.

\textit{Persistence.} Persistence assumes that opinions do not change and simply copies the
current opinion,
\[
\hat{s}_i(t{+}1)=s_i(t).
\]
Because opinion trajectories are often highly persistent, this baseline can be surprisingly
strong under plain accuracy despite predicting no opinion reversals. For the rollouts, persistence predicts $\hat{s}_i(t)=s_i(0)$.


\textit{Interaction-Free Model.} This baseline retains agent- and question-specific information while removing all social interactions. Predictions depend only on the intrinsic field,
\[
P\big(s_i(t{+}1)=+1\big)
=
\sigma\!\big(g_i\big),
\qquad
g_i={w}_{\text{field}}^{\top}{\phi}_i .
\]
Since the field is independent of the current opinion state, predictions are identical across
rounds and exhibit no dynamics.The model has
$(1+d_p+d_q+d_p d_q)$ learned parameters,\footnote{$d_p = d_q = 3$ and $1+d_p+d_q+d_p d_q = 16$ in all experiments.} the bias, persona and question components,
and their interaction. It isolates how much of an agent's next opinion is predictable
from who it is and what it was asked, before any communication with other agents.


\textit{Mean-Field Model.} This baseline incorporates social influence while discarding network
structure. Instead of interacting through the signed graph, each agent responds only to the
population-average opinion,
\vspace{-4pt}
\[
\bar{s}(t)=\frac{1}{n}\sum_j s_j(t),
\quad \text{yielding} \quad
P\big(s_i(t{+}1)=+1\big)
=
\sigma\!\big(g_i+\beta\,\bar{s}(t)\big).
\]
The model captures a conformity pressure toward the prevailing opinion while ignoring
heterogeneity in social ties. It assumes a $J$ in which every pair of agents is coupled with equal strength. Therefore, this baseline tests whether the signed-network structure provides additional predictive value. Relative
to the Interaction-Free baseline it introduces a single additional coupling parameter $\beta$, for a total of $(1+d_p+d_q+d_p d_q)+1 = 17$ learned parameters.
